\section{The imaginary unit and the Dirac structure}
\label{sec:dirac}

\subsection{The complex structure is selected}

At $n$ constituents write $\Omega=\gamma^0\gamma^1\gamma^2\gamma^3=e^{(1)}_1e^{(1)}_4e^{(1)}_5e^{(1)}_6$ for the volume element of the
interior four-plane and $\wI$ for the product, in increasing order, of the remaining $6n-4$ generators, the \emph{internal}
directions. At $n=1$ these are $e_2,e_3$ and $\wI=e_2e_3=J$.

\begin{theorem}[Selection of the imaginary unit]
\label{thm:i}
Inside $\End_\R(S_n)$ the joint commutant of the interior Lorentz algebra and of all internal bivectors is
$\operatorname{span}\{\id,\Omega,\wI,\Omega\wI\}\cong\C\otimes_\R\C\cong\C\oplus\C$, at every rung. Its solutions of
$x^2=-\id$ are the four isolated points $\pm\Omega$, $\pm\wI$. Of these only $\pm\wI$ commute with the interior gamma
matrices, from which the kinetic operator is built. Hence the imaginary unit is $i=\pm\wI$, and the chirality is
$\gamma_5=\mp\Omega\wI=\mp\Gamma$ with
\[
\Gamma:=\wI\,\Omega,\qquad \Gamma^2=\id .
\]
\end{theorem}

\begin{proof}
By \eqref{eq:bladecomm}, commuting with the six interior bivectors and with every internal bivector is a parity condition on
each blade, and the blades that satisfy all of them are exactly the four listed, at every rung. A four-dimensional
commutative real algebra with two commuting square roots of $-\id$ and no nilpotents is $\C\oplus\C$, whose square roots of
$-1$ are those four points. Finally $\Omega$ is the product of the four interior generators, so it anticommutes with each of
them, while $\wI$ contains no interior index and commutes with each.
\end{proof}

From now on $J$ denotes $\wI$; at one constituent this is $e_2e_3$, the generator of \cref{prop:6plus8plus1}. The selection
depends on its conditions, and it is worth saying exactly how.

\begin{remark}[The ambient algebra is part of the statement]
\label{rem:ambient}
If commutation with the internal bivectors is dropped, then at $n=1$ the commutant of the four gamma matrices in $M_8(\R)$ is the
quaternion algebra
\[
\HH=\operatorname{span}\{\id,\ J,\ e_2\Omega,\ e_3\Omega\},
\]
and its square roots of $-\id$ form a two-sphere. One further datum removes that freedom. The volume element of the temporal triple,
$T:=e_1e_2e_3$, depends only on the splitting $V=V^+\oplus V^-$ and on an orientation of $V^+$; it commutes with $J$ and anticommutes
with $e_2\Omega$ and $e_3\Omega$ by \eqref{eq:bladecomm}, so
\[
\{X\in\HH:\ [X,T]=0\}=\R\,\id\oplus\R\,J ,
\]
and $J=\gamma^0T$ is recovered from the clock and the temporal orientation, up to the sign of that orientation. Higher up the
tower the commutant of the gamma matrices alone acquires a matrix multiplicity algebra; at $n=3$ the Lorentz commutant is
$M_{128}(\C)$, whose own complex structure is $\Omega$ and not $\wI$ (\cref{prop:config}). Selection tracks the conditions imposed,
and every uniqueness statement in this paper is made with them.
\end{remark}

That the imaginary unit is not adjoined but read off the algebra is the cleanest structural result at one constituent. It has an
independent counterpart: on a real Hilbert space, Poincar\'e symmetry alone can force the complex structure \cite{MorettiOppio}.

\subsection{Real Dirac theory is complex Dirac theory}

\begin{theorem}[Exactness of the realification]
\label{thm:realification}
The $\C$-linear algebra generated by the interior gamma matrices on $(S,J)\cong\C^4$ is all of $M_4(\C)$. Consequently the real
Maxwell--Dirac system on $\Cl(3,3)$ is ordinary complex Maxwell--Dirac theory written in real variables.
\end{theorem}

\begin{proof}
$J$ commutes with the four $\gamma^\mu$. The sixteen monomials in the $\gamma^\mu$ and their $J$-multiples are thirty-two
linearly independent real matrices commuting with $J$, and the commutant of $J$ in $M_8(\R)$ is $\End_\C(\C^4)=M_4(\C)$, of real
dimension thirty-two.
\end{proof}

\begin{corollary}
\label{cor:nopredict}
Anything derivable from the real formulation alone was derivable from the complex one, and conversely. The gauged system is quantum
electrodynamics in real variables, and no prediction can follow from the rewriting.
\end{corollary}

This is the deflationary principle of the introduction. It does not make the arena uninteresting; it says where its content can
lie. Everything below that depends only on the four gamma matrices and $J$ is inherited. What is not inherited is what the arena
adds: the selection of $J$, the surplus plane and its uses, the pairing, and the tower.

\subsection{Admissible bilinears}

A field $\psi$ with values in $S$ is real. If its components anticommute, a term $\psi^{\!\top}A\psi$ vanishes when $A$ is
symmetric, and $\psi^{\!\top}A\,\partial\psi=\tfrac12\partial(\psi^{\!\top}A\psi)$ is a total derivative when $A$ is antisymmetric. A
mass term therefore needs an antisymmetric kernel, and a first-order kinetic term a symmetric one. With the clock $\gamma^0=e_1$:

\begin{proposition}[Kernels of fixed symmetry]
\label{prop:kernels1}
The mass form $\gamma^0J$ is antisymmetric. The kinetic kernels
\[
K^\mu:=(\gamma^0J)(J\gamma^\mu)=-\gamma^0\gamma^\mu,\qquad K^0=-\id,
\]
are symmetric, while $(\gamma^0J)\gamma^\mu$ is antisymmetric. The kinetic form must therefore carry the factor $J$: its $J$-free analogue
defines a total derivative and constrains nothing a field equation sees.
\end{proposition}

The mass form $\gamma^0J=e_1e_2e_3=T$ is the temporal volume element of \cref{rem:ambient}. It is not the only admissible mass.

\begin{proposition}[Scalar and pseudoscalar mass]
\label{prop:massangle}
The antisymmetric elements of $M_8(\R)$ that commute with $J$ and are invariant under the interior Lorentz algebra form the
two-dimensional space $\operatorname{span}\{T,\ T\Omega\}$. Used as the kernel of the Dirac bilinear, the second contributes nothing
to the kinetic term, because $T\Omega J\gamma^\mu$ is antisymmetric for every $\mu$; the two differ only in the mass term. $T$ is a
scalar mass and $T\Omega$ a pseudoscalar one, since $\Omega^2=-\id$ and $\Omega$ anticommutes with every $\gamma^\mu$, so that $\Omega$
is the real form of $i\gamma_5$. The family $T(\cos\vartheta+\sin\vartheta\,\Omega)=Te^{\vartheta\Omega}$ is the chiral mass angle.
\end{proposition}

So $T$ is singled out among admissible masses by the requirement that the mass be scalar, and that requirement is an input. The
central element $\Omega$ will reappear in \cref{sec:discrete} as the axial $U(1)$.

\paragraph{The Hamiltonian.} The free Dirac Hamiltonian at momentum $p$ is
\begin{equation}
\label{eq:Hp}
H(p)=\gamma^0\gamma^ap_a+m\gamma^0,\qquad H(p)^2=E^2\id,\qquad E^2=|\bp|^2+m^2 ,
\end{equation}
with spectral projectors $\Lambda_\pm(p)=\tfrac12(\id\pm H(p)/E)$. Both $\gamma^0=e_1$ and $\gamma^0\gamma^a=e_1e_{a+3}$ are symmetric, so
$H(p)$ is a symmetric matrix that commutes with $J$; its eigenvalues are $\pm E$, each with real multiplicity four.

\begin{example}
\label{ex:Hp}
At $\bp=(0,0,3)$ and $m=4$, $H=3e_1e_6+4e_1$ satisfies $H^2=(9+16)\id=25\,\id$, because $e_1e_6$ and $e_1$ anticommute and
$(e_1e_6)^2=+\id$. The eigenvalues are $\pm5$; each eigenspace is a complex plane for $J$, and the helicity
$\Sigma_3=J\,\tfrac12e_4e_5$ commutes with $H$ and splits each of them into two real planes.
\end{example}

\subsection{A first-order real form of free evolution}

For a commuting real field the same data give a first-order symplectic system. Take
\[
\cL_1=\tfrac12\Psi^{\!\top}J\dot\Psi-\tfrac12\Psi^{\!\top}H(p)\Psi .
\]
Its Euler--Lagrange equation is $\dot\Psi=-JH(p)\Psi$. Since $JH$ is antisymmetric with $(JH)^2=-E^2\id$, the flow lies in $SO(8)$ and
\[
\Psi(\lambda)=\cos(E\lambda)\,\Psi_0-\frac{\sin(E\lambda)}{E}\,JH\Psi_0 .
\]
Reading $J$ as $i$ turns this into $i\dot\psi=h\psi$ with $h=\alpha\cdot\bp+\beta m$, the free Dirac Hamiltonian at fixed momentum: an exact
embedding of free Dirac evolution in a real symplectic form, with eight second-class constraints reducing a sixteen-dimensional phase
space to eight. Three things norm conservation does not buy. The energy $\tfrac12\Psi^{\!\top}H\Psi$ is unbounded below, since $\Tr H=0$
and $H^2=E^2\id$ force the spectrum $\{\pm E\}$ with equal multiplicity. For anticommuting components the action collapses to a boundary
term, $\Psi^{\!\top}H\Psi=0$ because $H$ is symmetric and $\Psi^{\!\top}J\dot\Psi=\tfrac12\tfrac{d}{d\lambda}(\Psi^{\!\top}J\Psi)$ because $J$ is
antisymmetric, which is why a fermionic Lagrangian must use the kernels of \cref{prop:kernels1}. And the bilinear $\Psi^{\!\top}J\dot\Psi$
is not preserved by the boosts. The real form of Dirac theory in which a bivector commuting with the gamma matrices plays the role of $i$
is old \cite{Hestenes67}; what is specific here is that the bivector is the surplus rotation $J$, selected by \cref{thm:i}.
